% Track p1 -- The compact families at the third rung.
% Verification: investigations/photon_selection/verify/p1_families.py (claims p1-1 ... p1-47).
% New references: sections/bib_p1.tex.
\section{The compact families at the third rung}
\label{p1:sec:main}

The second-rung vacuum of \cite[Def.~26.5]{ES17} lies on a compact family $S^4$ of Lorentz-invariant projectors, and a spectrum-only
potential is constant on it \cite[Section~26]{ES17}. A coupling that distinguishes its points has to be found at $n\ge3$, where colour survives
\cite[Section~31]{ES17}. This section fixes the geometry on which any such coupling acts. The results are as follows. At the third rung the
Lorentz-invariant symmetric involutions that are products of four internal generators form three orbits of the compact group $K$. Only the
four-positive orbit $F_3$, an oriented Grassmannian of dimension $16$, commutes with colour, and it commutes with colour at every point
(fixed). The second-rung family is one Clifford four-sphere inside $F_3$ and is not $K$-invariant (fixed). The vacuum $\Pi_A$ built on the
auxiliary four-plane is the analogue of \cite[Def.~26.5]{ES17}. It is neutral under the construction's charge and under electromagnetism, and
it is fixed by the Pati--Salam and Standard-Model algebras (fixed). Its sixteen fluctuations are four complex Standard-Model doublets of
hypercharge $\tfrac12$ (fixed). Any $K$-invariant coupling is constant on $F_3$. The alignment energy of a proper gauged subalgebra vanishes
exactly on its fixed set, and these fixed sets are determined here. A $K$-equivariant duality $P\mapsto P^\perp$ preserves every gauge-boson
mass matrix, so no alignment energy separates $\Pi_A$ from $\Pi_W$ (fixed). On the matter block the construction's charge acts as half the
fermion number (fixed), and with $\su2_L$ it is anomalous there (restated). Identifiers p1-$n$ refer to the \texttt{PASS} lines of
\texttt{verify/p1\_families.py} (\cref{p1:tab:claims}).

\paragraph{Conventions.} The realisation is that of \cite[App.~A]{ES17}: $E_{6(k-1)+a}=\omega^{\otimes(k-1)}\otimes e_a\otimes\id^{\otimes(n-k)}$ on
$\R^{8^n}$, every generator and blade a signed permutation matrix, and $e_a^{\!\top}=\eta_{aa}e_a$. Products, commutators and traces of blades are therefore integer
computations, and blades obey the commutation rule, equation~(3) of \cite{ES17}, and the transpose rule, equation~(4) there.

\paragraph{Definitions.} Rung $n=3$, spectator interior action: interior generators $\gamma^\mu=(E_1,E_4,E_5,E_6)$; internal generators the other fourteen. The
positive internal span is $V_+=\operatorname{span}(A\cup W)$, of dimension $8$, and the negative internal span is $V_-=\operatorname{span}(C)$, of dimension $6$. The groups are those of \cite[Section~8]{ES17}: auxiliary $A=(E_2,E_3,E_8,E_9)$,
weak $W=(E_7,E_{13},E_{14},E_{15})$, written $(w_1,\dots,w_4)$, and colour $C=(E_{10},E_{11},E_{12},E_{16},E_{17},E_{18})$.
$K=\mathrm{Spin}(8)\times\mathrm{Spin}(6)$ is the group generated by the $43$ compact internal bivectors \cite[Thm.~2.3]{ES17}, with Lie algebra $\mathfrak k=\so8\oplus\so6$. For an
oriented four-plane $P$ in $V_+$ with oriented orthonormal basis $(u_1,u_2,u_3,u_4)$, $\mathrm{vol}_P:=u_1u_2u_3u_4$, a symmetric involution commuting with every interior
gamma, and $\Pi_P:=\tfrac12(\id+\mathrm{vol}_P)$. The compact family $F_3:=\{\Pi_P:P\text{ an oriented four-plane in }V_+\}$ is the $K$-orbit of $\Pi_A$ ($A$ the
auxiliary four-plane) and is diffeomorphic to the oriented Grassmannian of four-planes in $\R^8$, of dimension $16$ (\cref{p1:thm:classify}). The gauged subalgebras considered
are: $\mathfrak h_{43}=\mathfrak k$, all compact internal bivectors, as in \cite[Prop.~21.6]{ES17}; $\mathfrak h_{\rm PS}=\so4_W\oplus\so6_C$, the Pati--Salam algebra
$\su2_L\oplus\su2_R\oplus\su4$; $\mathfrak h_{\rm SM}=\su3_C\oplus\su2_L\oplus\uu1_Y$, embedded as in \cite[Section~8]{ES17}; $\uu1_{\rm em}$, the bivector part of
$Q=-JT^3_L+Y$; and $\uu1_A=\operatorname{span}(q_A)$, the construction's charge $Q=\tfrac12J_1$ at $n=3$. The alignment energy of a gauged subalgebra $\mathfrak h$ with
orthonormal basis $(T_a)$ and couplings $g_a$ is $V_{\mathfrak h}(\Pi):=\sum_ag_a^2\|[T_a,\Pi]\|^2$ (Frobenius norm on $\R^{512}$). With the projector kinetic term
$\tfrac{f^2}2\Tr(D\Pi D\Pi)$ of \cite[Section~26]{ES17} it is proportional to the trace of the gauge-boson mass matrix at $\Pi$.

\paragraph{Notation.} $J$ is the ordered product of the internal generators, $\Gamma=J\gamma^0\gamma^1\gamma^2\gamma^3$, $q_A=E_2E_3$, $Q_A=\tfrac12q_A$, $\chi_W=w_1w_2w_3w_4$,
$\chi_C$ the product of $C$, and $P_*=\tfrac18(\id+X)(\id+Z)(\id+U)$ with $X=\Gamma\chi_W$, $Z=-J\chi_C$, $U=-Jq_A$ \cite[Thm.~8.1]{ES17}; $P_{L,R}=P_*\tfrac12(\id\mp\Gamma)$.
The colour Cartan element is $H_C=-(E_{10}E_{16}+E_{11}E_{17}+E_{12}E_{18})$, and $\su3_C$ is its derived centraliser in $\so6_C$ \cite[Prop.~8.4]{ES17}. Further
$\tau_i=-\tfrac12w_jw_k$ for cyclic $(i,j,k)$, $t^{L,R}_i=\tfrac12(\id\mp\chi_W)\tau_i$ \cite[Thm.~8.2]{ES17}, the hypercharge bivector is $t_Y=t^R_3+\tfrac16H_C$, and
\[
B_{\rm em}=\tau_3+\tfrac16H_C=t^L_3+t_Y=-\tfrac12E_7E_{13}+\tfrac16H_C
\]
spans $\uu1_{\rm em}$: on the block $Q=-J(t^3_L+t^3_R+\tfrac16H_C)$ \cite[Prop.~8.4]{ES17}. The vector representation $\rho$ is defined by $[X,u]=\rho(X)u$ for internal
vectors $u$. Thus $\rho(\tfrac12e_ie_j)=E_{ij}-E_{ji}$ on $V_+$, with a sign $-1$ on $V_-$, and the lift $M\mapsto\sum_{i<j}M_{ij}\tfrac12e_ie_j$ (sign $-1$ on $V_-$) is a Lie
algebra map inverse to $\rho$ (p1-12). We write $\rho_+(X)\in\so8$ for the $V_+$ part, $p_P$ for the orthogonal projector of $\R^8$ onto $P$, $\theta$ for $+1$ on $V_+$ and
$-1$ on $V_-$, and $\widetilde{\mathrm{Gr}}_k(V)$ for the oriented Grassmannian. On vectors, $\rho(Q_A)$ maps $E_2\mapsto-E_3$, $E_3\mapsto E_2$ and kills $E_8,E_9,W$. $\rho(B_{\rm em})$
maps $w_1\mapsto w_2$, $w_2\mapsto-w_1$ and kills $A,w_3,w_4$. $J_Y=2\rho(t_Y)$ is a complex structure on $W$ that kills $A$.

\subsection{The second rung, restated}
\label{p1:sec:n2}

\begin{proposition}[The second-rung family is a $K_2$-orbit; restated and fixed]
\label{p1:prop:n2}
At $n=2$ the internal complement of the spectator action has signature $(5,3)$, with $V_+^{(2)}=\operatorname{span}(E_2,E_3,E_7,E_8,E_9)$ and
$V_-^{(2)}=\operatorname{span}(E_{10},E_{11},E_{12})$. Let $(a,b,c,d,e)=(E_2,E_3,E_7,E_8,E_9)$, $Q=\tfrac12ab$, $H=abcd$, $\Pi_0=\tfrac12(\id+H)$, $I_5=abcde$, $Z_i=I_5v_i$.
\begin{enumerate}[label=(\roman*),leftmargin=2.2em]
\item \emph{(Restated.)} \cite[Thms.~26.6--26.8]{ES17} hold: $H$ is a symmetric traceless involution commuting with the interior gammas, $Q$ and the
internal volume; $[\tfrac12be,\Pi_0]=-\tfrac12acde$; the $Z_i$ are a symmetric Clifford system, $T_i=[\tfrac12v_ie,\Pi_0]=\tfrac12Z_i$ with $\Tr(T_iT_j)=16\delta_{ij}$, $Q$
rotates $(T_a,T_b)$ and fixes $T_c,T_d$; along $\Pi(n)=\tfrac12(\id+I_5n)$, $\Tr(d\Pi^2)=16|dn|^2$, $-\Tr([Q,\Pi]d\Pi)=16(n_adn_b-n_bdn_a)$ and $[\Pi,d\Pi]=\tfrac14[n,dn]$.
The stabiliser of $H$ in $\so{5,3}$ is $\so P\oplus\so{P^\perp}=\so4\oplus\so{1,3}$, of dimension $12$ with compact part $9$, and every other internal bivector
anticommutes with $H$ \cite[Prop.~26.9]{ES17}.
\item \emph{(Fixed.)} The $13$ compact internal bivectors preserve $\operatorname{span}\{Z_i\}$: $\so5$ acts on it by the vector representation and $\so3$
trivially. The $K_2$-orbit of $\Pi_0$, $K_2=\mathrm{Spin}(5)\times\mathrm{Spin}(3)$, is the whole family $S^4=\{\Pi(n)\}$, and
$\exp(\tfrac\vartheta2v_ie)$ carries $\Pi(e)$ to $\Pi(\cos\vartheta\,e+\sin\vartheta\,v_i)$. Moreover $I_5n=\mathrm{vol}_{n^\perp}$, so $S^4=\{\Pi_P:P\in\widetilde{\mathrm{Gr}}_4(V_+^{(2)})\}$.
\item \emph{(Fixed.)} $[Q,\Pi(n)]=0$ exactly on the two-sphere $n_a=n_b=0$, the planes containing $\operatorname{span}(a,b)$. Every $\Pi(n)$ has spectrum
$\{0^{32},1^{32}\}$, and $\tfrac1{64}\sum\|[e_ie_j,\Pi(n)]\|^2=4$ over the $13$ compact blades. A spectrum-only potential and the full alignment energy are
constant on $S^4$.
\end{enumerate}
\end{proposition}

\begin{proof}
(i) is the computation of \cite[Section~26]{ES17}, repeated in p1-1 to p1-4 and p1-7. The stabiliser statement is equation~(3) of
\cite{ES17}: a bivector $e_ie_j$ and the grade-four blade $H$ commute iff $|\{i,j\}\cap\{a,b,c,d\}|$ is even. (ii) A positive bivector inside $V_+^{(2)}$ shares two
indices with $I_5$ and commutes with it, so $[X,Z_i]=I_5[X,v_i]=I_5\rho(X)v_i$; a negative bivector shares none and commutes with each $Z_i$. Since $\mathrm{Spin}(5)$ is
transitive on $S^4$, the orbit of $\Pi(e)$ is the family. The rotation formula is $\Ad(\exp\tfrac\vartheta2v_ie)e=\cos\vartheta\,e+\sin\vartheta\,v_i$.
If $(f_1,\dots,f_4,n)$ is an oriented orthonormal basis of $V_+^{(2)}$, then $I_5=f_1f_2f_3f_4n$ and $I_5n=f_1f_2f_3f_4$ (p1-5). (iii) $[Q,Z_i]=I_5\rho(Q)v_i$, and
$\rho(Q)$ rotates $(a,b)$ and kills $c,d,e$ (p1-6). Each $\Pi(n)$ is conjugate to $\Pi_0$ by an orthogonal matrix, and the alignment sum is $K_2$-invariant, as in
\cref{p1:prop:K} below (p1-4, p1-5).
\end{proof}

At $n=2$ three descriptions coincide: the Clifford sphere of five positive generators, the $K_2$-orbit, and the oriented Grassmannian of four-planes in
$V_+^{(2)}$. At $n=3$ they separate. The Grassmannian description is the one that continues.

\subsection{Four-generator involutions at the third rung}
\label{p1:sec:classify}

\begin{lemma}[Blade census; fixed]
\label{p1:lem:census}
At $n=3$ the eighteen generators satisfy the Clifford relations of $\R^{9,9}$, $V_+=\operatorname{span}(E_2,E_3,E_7,E_8,E_9,E_{13},E_{14},E_{15})$, and the $43$ compact
internal bivectors square to $-\id$ while the $48$ mixed ones square to $+\id$ (p1-8).
\begin{enumerate}[label=(\roman*),leftmargin=2.2em]
\item Of the $3060$ blades of grade four exactly the $1001$ internal ones commute with the four interior gammas. All are traceless. One with $q$ negative factors is
a symmetric involution if $q$ is even and an antisymmetric complex structure if $q$ is odd. The counts are $(4,0):70$, $(3,1):336$, $(2,2):420$, $(1,3):160$, $(0,4):15$
(p1-9).
\item More generally an internal blade of even grade $r$ with $q$ negative factors is a symmetric involution iff $r(r-1)/2+q$ is even. The only nontrivial one invariant
under $K$ is $\mathrm{vol}_+=E_2E_3E_7E_8E_9E_{13}E_{14}E_{15}$, and $\mathrm{vol}_+=-Z=J\chi_C$ (p1-10).
\end{enumerate}
\end{lemma}

\begin{proof}
For $|I|=4$, equation~(3) of \cite{ES17} gives $e_Ie_j=(-1)^{4-|I\cap\{j\}|}e_je_I$, so a generator commutes with $e_I$ iff it is not a factor. By equation~(4) of
\cite{ES17}, $e_I^2=(-1)^{r(r-1)/2+q}\id$ and $e_I^{\!\top}=e_I^2e_I$. An even internal blade commutes with $e_ie_j$ iff $|I\cap\{i,j\}|$ is even. For all same-sign pairs this
forces $I\cap V_\pm\in\{\emptyset,V_\pm\}$, which leaves $\mathrm{vol}_+$, $\chi_C$ ($r=q=6$, square $-\id$) and $\pm J$ ($r=14$, square $-\id$). The
identity $\mathrm{vol}_+=J\chi_C$ is checked on the matrices.
\end{proof}

\begin{lemma}[Transpose and recovery; fixed]
\label{p1:lem:transpose}
Let $P$ be a nondegenerate internal four-plane with $\eta$-orthonormal oriented basis $(u_1,\dots,u_4)$ and $\mathrm{vol}_P=u_1u_2u_3u_4$. Then $\mathrm{vol}_P$ commutes with the
interior gammas, $\mathrm{vol}_P^2=(-1)^q\id$, the internal vectors anticommuting with $\mathrm{vol}_P$ are exactly those of $P$, and
\[
\mathrm{vol}_P^{\!\top}=\mathrm{vol}_{\theta P}.
\]
Hence for $q$ even, $\Pi_P$ is an orthogonal projector iff $\theta P=P$, that is iff $P=(P\cap V_+)\oplus(P\cap V_-)$. For a positive-definite $P\not\subset V_+$, $\Pi_P$ is
an idempotent but not symmetric (p1-13).
\end{lemma}

\begin{proof}
The interior gammas anticommute with every internal vector, hence commute with a product of four. By equation~(4) of \cite{ES17}, $\mathrm{vol}_P^2=(-1)^6\prod\eta(u_i)\id=(-1)^q\id$.
For an internal vector $u$, $u^{\!\top}=\theta u$. Transposing reverses the product, and reversing four anticommuting vectors costs $(-1)^{6}=1$. So
$\mathrm{vol}_P^{\!\top}=\theta u_1\cdots\theta u_4=\mathrm{vol}_{\theta P}$, and $\theta$ is an isometry of $\eta$. Each $u_i$ commutes with itself and anticommutes with the other three
factors, so every vector of $P$ anticommutes with $\mathrm{vol}_P$. A vector $\eta$-orthogonal to $P$ anticommutes with all four factors, so it commutes with $\mathrm{vol}_P$. Since $P$ is
nondegenerate, an internal $v$ splits as $v_P+v_\perp$; if $v$ anticommutes with $\mathrm{vol}_P$, so does $v_\perp$, which also commutes, hence $v_\perp\mathrm{vol}_P=0$ and $v_\perp=0$.
So $P$ is determined by $\mathrm{vol}_P$, and its orientation by the sign. Then $\mathrm{vol}_P$ is symmetric iff $\theta P=P$ as oriented planes.
\end{proof}

\begin{lemma}[Colour acts irreducibly on $V_-$; fixed]
\label{p1:lem:colour}
The centraliser of $\rho(H_C)$ in $\so{V_-}$ is $\uu3$. Its part orthogonal to $H_C$ is $\su3_C$, of dimension $8$ and perfect. Lifted to $\mathfrak k$, $\su3_C$ commutes
with $H_C$ and acts trivially on $V_+$. Its commutant on $V_-=\R^6$ is $\operatorname{span}(\id,\tfrac12\rho(H_C))$ with $(\tfrac12\rho(H_C))^2=-\id$, a field, so
$\su3_C$ acts irreducibly on $V_-$ (p1-15).
\end{lemma}

\begin{proof}
The computation is in the six-dimensional vector model and is lifted by the Lie algebra map of p1-12. A proper invariant subspace of an orthogonal representation
has an invariant orthogonal complement, and the projector onto it is an idempotent in the commutant other than $0$ and $\id$. A field has no such idempotent.
\end{proof}

\begin{theorem}[The three orbit types and the colour test; fixed]
\label{p1:thm:classify}
The Lorentz-invariant symmetric involutions $\mathrm{vol}_P$, with $P$ a nondegenerate internal four-plane, form three $K$-orbits. Their stabilisers and dimensions are:
\begin{center}
\footnotesize
\begin{tabular}{@{}lllccc@{}}
\toprule
type & orbit & stabiliser in $\so{8,6}$ & dim & compact part & orbit dim\\
\midrule
$(4,0)$ & $F_3\cong\widetilde{\mathrm{Gr}}_4(V_+)$ & $\so4\oplus\so{4,6}$ & $51$ & $27$: $\so4_A\oplus\so4_W\oplus\so6_C$ & $16$\\
$(2,2)$ & $\widetilde{\mathrm{Gr}}_2(V_+)\times\widetilde{\mathrm{Gr}}_2(V_-)/\pm$ & $\so{2,2}\oplus\so{6,4}$ & $51$ & $23$ & $20$\\
$(0,4)$ & $\widetilde{\mathrm{Gr}}_4(V_-)$ & $\so4\oplus\so{8,2}$ & $51$ & $35$ & $8$\\
\bottomrule
\end{tabular}
\end{center}
Every point has spectrum $\{+1^{256},-1^{256}\}$, so no spectrum-only potential separates the points or the types. Every point of $F_3$ commutes with all of
$\so6_C\supset\su3_C$. No point of the $(2,2)$ or $(0,4)$ orbits commutes with $\su3_C$.
\end{theorem}

\begin{proof}
By \cref{p1:lem:transpose} the symmetric ones are the $\mathrm{vol}_{P_+}\mathrm{vol}_{P_-}$ with $P_\pm\subset V_\pm$ of dimensions $(p,q)$, $q$ even. For $g\in K$,
$\Ad(g)u=\rho(g)u$ on vectors, so $\Ad(g)\mathrm{vol}_P=\mathrm{vol}_{\rho(g)P}$ (p1-14). $SO(8)\times SO(6)$ is transitive on pairs of oriented planes, which gives one orbit per
type; reversing both orientations in type $(2,2)$ gives the same element. By the recovery statement of \cref{p1:lem:transpose}, $\exp(tX)$ fixes $\mathrm{vol}_P$ iff $\rho(X)$
preserves $P$, so the stabiliser is $\so P\oplus\so{P^\perp}$. Bivector by bivector, $e_ie_j$ commutes with $\mathrm{vol}_P$ iff both or neither of $i,j$ lie in $P$, and
anticommutes otherwise; the counts are p1-11. The orbit dimension is $43$ minus the compact part. A symmetric involution of trace zero has eigenvalues $\pm1$, each $256$
times (p1-42). A vector of $V_+$ anticommutes with both factors of a colour bivector and so commutes with it; hence every colour bivector commutes with every $\mathrm{vol}_P$, $P\subset V_+$. For types $(2,2)$ and $(0,4)$,
$\su3_C$ fixes $\mathrm{vol}_P$ iff it preserves the plane $P_-\subset V_-$ of dimension $2$ or $4$, which \cref{p1:lem:colour} excludes. p1-16 checks this at the
representatives and at random $K$-translates.
\end{proof}

At $n=3$ the stabiliser of the four-positive type is the row of \cite[Prop.~26.9]{ES17}. The same proof covers all even grades: a $\theta$-stable plane of type $(p,q)$
with $r(r-1)/2+q$ even gives a $K$-orbit of dimension $p(8-p)+q(6-q)$. Among them only $\mathrm{vol}_+=-Z$ is $K$-fixed (\cref{p1:lem:census}). No compact generator moves
$\tfrac12(\id+\mathrm{vol}_+)$, as for the chirality projectors of \cite[Prop.~26.4]{ES17}; only the boosts do. \Cref{p1:lem:transpose} agrees with \cite[Thm.~11.3]{ES17}, by which a
variation $[B,\Pi]$ stays among the orthogonal projectors exactly when $B$ is compact: a noncompact generator tilts $P$ out of $V_+$, and $\Pi_P$ stops being symmetric.

\begin{proposition}[Clifford spheres in $F_3$; fixed]
\label{p1:prop:spheres}
\begin{enumerate}[label=(\roman*),leftmargin=2.2em]
\item For an oriented five-plane $L\subset V_+$ with oriented orthonormal basis $(v_1,\dots,v_5)$ put $I_L=v_1\cdots v_5$ and $Z_i=I_Lv_i$. The $Z_i$ are a symmetric Clifford
system, and $\tfrac12(\id+I_Ln)=\Pi_{n^\perp\cap L}$. Hence $S_L:=\{\tfrac12(\id+I_Ln):|n|=1\}=\{\Pi_P:P\subset L\}\cong S^4$ is the $\mathrm{Spin}(L)$-orbit of any of its points.
Dually, for an oriented three-plane $U$ with volume $U_3$, $S^U:=\{\Pi_P:P\supset U\}=\{\tfrac12(\id+nU_3):n\in U^\perp,|n|=1\}\cong S^4$.
\item Through each $\Pi_P$ pass the spheres $S_L$ with $L=P\oplus\R u$, $u\in P^\perp$ (a projective space $\R P^3$ of them), with tangent $P\otimes u$, and the spheres $S^U$
with $U=P\cap a^\perp$, $a\in P$, with tangent $a\otimes P^\perp$. These are the two kinds of maximal linear spaces of rank-one elements in $T_{\Pi_P}F_3\cong P\otimes P^\perp$. The
inner spheres are parametrised by $\mathrm{Gr}_5(\R^8)$, of dimension $15$, and none of them is $K$-invariant.
\item On the inner sphere through $A\oplus\R E_7$ with $Q=Q_A$, and on the outer sphere of planes containing $\operatorname{span}(E_3,E_8,E_9)$ with $Q=B_{\rm em}$,
$\Tr(d\Pi^2)=128|dn|^2$ and $-\Tr([Q,\Pi]d\Pi)=\pm128(n_adn_b-n_bdn_a)$. On the inner sphere $[\Pi,d\Pi]=\tfrac14[n,dn]$. This is the structure of
\cite[Thms.~26.7, 26.8]{ES17} with $16$ replaced by $128$ (p1-19, p1-32).
\item The second-rung family embedded with its own generators, $L_0=\operatorname{span}(a,b,c,d,e)=\operatorname{span}(E_2,E_3,E_7,E_8,E_9)$, is the inner sphere $S_{L_0}$. Its
point $n=c$ is $\Pi_A$, since $I_5c=E_2E_3E_8E_9$, and its point $n=e$ is $\Pi_0=\tfrac12(\id+E_2E_3E_7E_8)$. On $S_{L_0}$, $Q_A=\tfrac12ab$ is unbroken on the two-sphere
$n_a=n_b=0$. The algebras $\uu1_{\rm em}$, $\mathfrak h_{\rm SM}$ and $\mathfrak h_{\rm PS}$ are each unbroken only at $n=\pm c$, that is at $\Pi_{\pm A}$ (p1-47).
\end{enumerate}
\end{proposition}

\begin{proof}
(i) As in \cref{p1:prop:n2}(ii): $Z_iZ_j+Z_jZ_i=2\delta_{ij}$ because $I_L$ commutes with each $v_i$ and squares to $+\id$, and $I_Ln=f_1f_2f_3f_4$ for an oriented basis
$(f_1,\dots,f_4,n)$ of $L$. For $S^U$, $nU_3$ is the volume of $\operatorname{span}(n,U)$; vectors of $U^\perp$ anticommute with $U_3$, so the $nU_3$ again form a Clifford system
(p1-19). (ii) The orbit of $\mathrm{Spin}(L)$, respectively $\mathrm{Spin}(U^\perp)$, has tangent $[\mathfrak m_P\cap\so L,\Pi_P]$, respectively $[\mathfrak m_P\cap\so{U^\perp},\Pi_P]$,
where $\mathfrak m_P=P\wedge P^\perp$ is spanned by the bivectors mixing $P$ and $P^\perp$. These intersections are $P\wedge u$ and $a\wedge P^\perp$. (iii) $\Pi$ and $d\Pi$ are
combinations of the $Z_i$, and $\Tr(Z_iZ_j)=512\delta_{ij}$; $[Q,Z_i]=\pm Z_j$ by the action of $\rho(Q)$ on $L$ or on $U^\perp$ (p1-32). (iv) The identities are blade
identities. The fixed points follow from \cref{p1:thm:fixed} below, since $P\subset L_0$ contains or is orthogonal to $\operatorname{span}(E_7,E_{13})$ only for $P=A$, and
$E_{13}\notin L_0$; p1-47 computes the kernels of $n\mapsto[X,I_5n]$ directly.
\end{proof}

\begin{proposition}[The duality $P\mapsto P^\perp$; fixed]
\label{p1:prop:duality}
$\mathrm{vol}_+$ commutes with $\mathfrak k$, with the interior gammas and with every $\mathrm{vol}_P$, and $\mathrm{vol}_+\mathrm{vol}_A=\mathrm{vol}_W$. The map
$D(\Pi)=\tfrac12\bigl(\id+\mathrm{vol}_+(2\Pi-\id)\bigr)$ sends $\Pi_P$ to $\Pi_{P^\perp}$, with $P^\perp\subset V_+$ suitably oriented. It is a $K$-equivariant involution of $F_3$
without fixed points, and $D(\Pi_A)=\Pi_W$. For every $T,T'\in\mathfrak k$,
\[
\Tr\bigl([T,D\Pi][T',D\Pi]\bigr)=\Tr\bigl([T,\Pi][T',\Pi]\bigr),\qquad \Tr\bigl([T,\id-\Pi][T',\id-\Pi]\bigr)=\Tr\bigl([T,\Pi][T',\Pi]\bigr).
\]
So every gauge-boson mass matrix, and every alignment energy $V_{\mathfrak h}$ with $\mathfrak h\subset\mathfrak k$, takes equal values at $\Pi_{\pm A}$ and $\Pi_{\pm W}$. On the
block, $\mathrm{vol}_+P_*=-P_*$, and $\Tr(\Pi_AP_L)=32$, $\Tr(\Pi_WP_L)=0$, $\Tr(\Pi_WP_R)=32$ (p1-20, p1-46).
\end{proposition}

\begin{proof}
$\mathrm{vol}_+$ is $K$-invariant by \cref{p1:lem:census}, and it commutes with the interior gammas, with which it shares no index. In an orthonormal basis $(u_1,\dots,u_8)$
of $V_+$ adapted to $P\oplus P^\perp$, $\mathrm{vol}_+=\pm u_1\cdots u_8$, so $\mathrm{vol}_+$ commutes with $\mathrm{vol}_P=u_1u_2u_3u_4$ and
$\mathrm{vol}_+\mathrm{vol}_P=\pm u_5u_6u_7u_8=\pm\mathrm{vol}_{P^\perp}$. $D$ is an involution because $\mathrm{vol}_+^2=\id$, and it has no fixed point because $P\neq P^\perp$. Since $\mathrm{vol}_+$
commutes with $T$ and $\Pi$, $[T,D\Pi]=\mathrm{vol}_+[T,\Pi]$, and $\mathrm{vol}_+^2=\id$ gives the identity. The second identity is $[T,\id-\Pi]=-[T,\Pi]$. On the block,
$\mathrm{vol}_+P_*=-ZP_*=-P_*$. The traces follow from \cref{p1:prop:block}.
\end{proof}

\subsection{The vacuum \texorpdfstring{$\Pi_A$}{Pi\_A}}
\label{p1:sec:vacuum}

\begin{definition}[The third-rung vacuum]
\label{p1:def:vacuum}
$\Pi_A:=\tfrac12(\id+E_2E_3E_8E_9)$, built from the surplus pairs of the first and second constituents, and $\Pi_W:=\tfrac12(\id+\chi_W)$. Their orientation
reversals are $\Pi_{-A}=\id-\Pi_A$ and $\Pi_{-W}=\id-\Pi_W$.
\end{definition}

\begin{theorem}[Lorentz invariant, neutral and not rigid; fixed]
\label{p1:thm:vacuum}
$\Pi_A$ and $\Pi_W$ are rank-$256$ orthogonal projectors with the following properties.
\begin{enumerate}[label=(\roman*),leftmargin=2.2em]
\item Each commutes with the interior gammas, with $J$, with $q_A$, with $B_{\rm em}$, with the $21$ generators of $\mathfrak h_{\rm PS}$ and with the $12$ of
$\mathfrak h_{\rm SM}$ (p1-21).
\item The compact internal bivectors that move $\Pi_A$ are exactly the $16$ that mix $A$ with $W$. The tangents $T_{aw}=[\tfrac12aw,\Pi_A]=\tfrac12aw\,\mathrm{vol}_A$ are
symmetric with $\Tr(T_{aw}T_{a'w'})=128\,\delta_{aa'}\delta_{ww'}$; the same holds at $\Pi_W$ (p1-23).
\item The stabiliser of $\Pi_A$ in $\mathfrak k$ is $\so4_A\oplus\so4_W\oplus\so6_C$, of dimension $27$. It contains $\uu1_A$, $\uu1_{\rm em}$, $\mathfrak h_{\rm PS}$ and
$\mathfrak h_{\rm SM}$ (p1-35).
\end{enumerate}
The literal embedding $\tfrac12(\id+E_2E_3E_7E_8)$ of the second-rung vacuum is also a point of $F_3$ and commutes with $q_A$. It does not commute with $B_{\rm em}$ or with
$\su2_L$, because $E_7\in W$ (p1-22).
\end{theorem}

\begin{proof}
All statements are equation~(3) of \cite{ES17} applied to blades. $\mathrm{vol}_A$ shares two indices with $q_A$, none with $E_7E_{13}$, the colour bivectors or the
weak bivectors, and four with $J$. A bivector $aw$ shares one index with $\mathrm{vol}_A$ and anticommutes with it, so
$[\tfrac12aw,\Pi_A]=\tfrac14[aw,\mathrm{vol}_A]=\tfrac12aw\,\mathrm{vol}_A$. Symmetry and the trace follow from $(aw\,\mathrm{vol}_A)^2=-(aw)^2=\id$ and the tracelessness of
nontrivial blades. $B_{\rm em}$, $t^{L,R}_i$ and $\su3_C$ are combinations of bivectors inside $W$ or $C$.
\end{proof}

In \cite[Prop.~10.12]{ES17} the compact spectator algebra carries colour on three negative planes, weak isospin on two positive planes and a family $\so4$ on the remaining
two. With the grouping of \cite[Section~8]{ES17} the remaining two planes are those of $A$, so $\Pi_A$ separates the weak four-plane from the family four-plane and keeps
everything else (interpretive). By \cref{p1:prop:spheres}(iv), $\Pi_A$ is the point $n=c$ of the second-rung family itself, embedded at $n=3$. The point $n=e$ chosen in
\cite[Def.~26.5]{ES17} is not neutral at $n=3$.

\begin{theorem}[The fluctuations at $\Pi_A$; fixed]
\label{p1:thm:tangent}
Let $T_{\Pi_A}F_3=\operatorname{span}\{T_{aw}\}\cong A\otimes W$.
\begin{enumerate}[label=(\roman*),leftmargin=2.2em]
\item Under $\so4_A\oplus\so4_W\cong\su2^{\oplus4}$ it is the real irreducible $(\mathbf2,\mathbf2;\mathbf2,\mathbf2)$, with commutant $\R$ and every $\su2$ Casimir equal to
$\tfrac34$ (p1-24).
\item Under $\mathfrak h_{\rm SM}$, $\su3_C$ acts trivially; $\su2_L$ has Casimir $\tfrac34$, and $\ad t_Y$ satisfies $(\ad t_Y)^2=-\tfrac14$ and commutes with $\su2_L$. The space is four
complex doublets $(\mathbf1,\mathbf2)_{1/2}$, one for each direction of $A$, and $\so4_A$ rotates the doublet index (p1-25).
\item Under $\uu1_{\rm em}$, $\ad B_{\rm em}=\ad(t^3_L+t_Y)$. The charged directions are $A\otimes\operatorname{span}(w_1,w_2)$, four complex scalars
$\phi_a=x_{a,13}+ix_{a,7}$ of charge one, the upper components of the doublets. The neutral directions are $A\otimes\operatorname{span}(w_3,w_4)$, eight real, the lower
components.
\item Under $\uu1_A$ the charged directions are $\operatorname{span}(E_2,E_3)\otimes W$, four complex scalars $\phi_w=x_{2,w}+ix_{3,w}$. The neutral directions are
$\operatorname{span}(E_8,E_9)\otimes W$, eight real.
\item The same holds at $\Pi_W$ with $\operatorname{Hom}(W,A)\cong W\otimes A$ (p1-24, p1-25, p1-26).
\end{enumerate}
Here $x_{a,w}$ are the coordinates of $\Pi=\Ad(e^X)\Pi_A$, $X=\sum x_{a,w}\tfrac12aw$. At $n=2$ the same count gives one charged complex scalar and two neutral real
scalars \cite[Thm.~26.7]{ES17}.
\end{theorem}

\begin{proof}
For $X$ in the stabiliser, $[X,[Y,\Pi_A]]=[[X,Y],\Pi_A]$. So $T_{\Pi_A}F_3$ is the representation of the stabiliser on $\mathfrak m_A=A\wedge W\cong A\otimes W$, on which
$\rho(\tfrac12aw)$ is the rotation of the $(a,w)$ plane. All five statements reduce to the action on $A$ and $W$ listed in the conventions. $\su2_L$ acts on $W$ with a
quaternionic commutant (p1-37), so $W$ is one complex doublet with complex structure $J_Y=2\rho(t_Y)$. $\rho(t^3_L+t_Y)=\rho(\tau_3)$ rotates $(w_1,w_2)$ at unit rate and
kills $(w_3,w_4)$. p1-26 gives the explicit pairs, for example $[Q_A,T_{2,w}]=-T_{3,w}$, $[Q_A,T_{3,w}]=T_{2,w}$ and $[B_{\rm em},T_{a,7}]=T_{a,13}$, $[B_{\rm em},T_{a,13}]=-T_{a,7}$.
\end{proof}

\begin{proposition}[The vacuum on the matter block; fixed, and restated]
\label{p1:prop:block}
\begin{enumerate}[label=(\roman*),leftmargin=2.2em]
\item $\mathrm{vol}_A$ and $\chi_W$ commute with $X,Z,U$, and $\mathrm{vol}_AP_*=-\Gamma P_*$, $\chi_WP_*=\Gamma P_*$. Hence $\Pi_AP_*=P_L$ and $\Pi_WP_*=P_R$, the two Weyl blocks of rank
$32$ (p1-27).
\item $q_AP_*=JP_*$. So $Q_A=\tfrac12q_A$, the third-rung form of the construction's charge $\tfrac12J_1$, acts on the block as $\tfrac12J$: its charge operator
$-JQ_A$ equals $\tfrac12$ on all sixteen Weyl components. It is half the fermion number of \cite[Section~8]{ES17}, not the electric charge, which takes the values
$-1,-\tfrac13,0,\tfrac23$ (p1-18, p1-27).
\item Every $T_{aw}$ anticommutes with $X=\Gamma\chi_W$, and $P_*T_{aw}P_*=0$: the fluctuations of $\Pi_A$ map the block into its orthogonal complement (p1-27).
\item \emph{(Restated, \cite[Prop.~8.5]{ES17} with $a=\tfrac12$, $b=0$.)} With $\hat Q_A=-JQ_A$ and $\hat T=-Jt^L_3$, we have $\Tr(P_L\hat Q_A\hat T^2)=4$ and
$\Tr(P_R\hat Q_A\hat T^2)=0$, and likewise $4$ for $\su2_R$ on $P_R$. For $\hat Q_{\rm em}=-JB_{\rm em}$, $\Tr(P_L\hat Q_{\rm em}\hat T^2)=0$. The mixed $SU(2)_L^2\,U(1)_A$ anomaly does
not cancel on the block, while the gravitational one does (p1-28) \cite{BouchiatIliopoulosMeyer72,GeorgiGlashow72}.
\end{enumerate}
\end{proposition}

\begin{proof}
(i) $\mathrm{vol}_A=\mathrm{vol}_+\chi_W$ by \cref{p1:prop:duality}. $\mathrm{vol}_+$ commutes with $\Gamma$ and $\mathrm{vol}_+P_*=-P_*$; with $\chi_WP_*=\Gamma P_*$
\cite[Thm.~8.1]{ES17} this gives $\mathrm{vol}_AP_*=-\Gamma P_*$. (ii) $UP_*=P_*$ means $-Jq_AP_*=P_*$, and $J^2=-\id$ gives $q_AP_*=JP_*$; then $-JQ_AP_*=\tfrac12P_*$. (iii)
$T_{aw}$ is half a blade of grade four with three factors in $A$ and one in $W$. It shares one index with $\chi_W$ and anticommutes with it, and it commutes with $\Gamma$,
which is $\pm$ the product of all eighteen generators. Hence it anticommutes with $X$ (p1-27).
Since $XP_*=P_*$, $P_*TP_*=P_*\tfrac12(\id+X)TP_*=P_*T\tfrac12(\id-X)P_*=0$. (iv) On $P_L$, $\hat Q_A=\tfrac12$ and $\hat T^2=-(t^L_3)^2=\tfrac14$, and $\Tr P_L=32$; $t^L_3P_R=0$. The electric charges of the left doublets sum to zero.
\end{proof}

The charge that the second-rung construction gauges, $Q=\tfrac12J_1$, therefore becomes at $n=3$ a universal charge on the
$P_*$ matter, and with $\su2_L$ it cannot be gauged on that matter. The charge for which $\Pi_A$ plays the role of \cite[Def.~26.5]{ES17} with electromagnetic content is
$B_{\rm em}$. Both meet the massless criterion of \cite[Cor.~26.2]{ES17} at $\Pi_A$ (\cref{p1:thm:kinetic}). Which of them is the photon is an identification and is
not decided by the geometry of $F_3$ (open).

\subsection{Kinetic term, current and transport on \texorpdfstring{$F_3$}{F\_3}}
\label{p1:sec:kinetic}

\begin{lemma}[Vector model; fixed]
\label{p1:lem:vector}
For $P\in\widetilde{\mathrm{Gr}}_4(V_+)$ and $X,Y\in\mathfrak k$,
\[
\Tr\bigl([X,\Pi_P][Y,\Pi_P]\bigr)=64\,\tr_8\bigl([\rho_+X,p_P][\rho_+Y,p_P]\bigr).
\]
Hence on $F_3$: $\tfrac{f^2}2\Tr(D\Pi D\Pi)=32f^2\tr_8(Dp\,Dp)$; for $Q\in\mathfrak k$, $j_\Pi=-f^2\Tr([Q,\Pi]D\Pi)=-64f^2\tr_8([\rho_+Q,p]Dp)$; and
\[
V_{\mathfrak h}(\Pi_P)=64\sum_ag_a^2\,\|[\rho_+T_a,p_P]\|^2 ,
\]
in which the $\so6$ components of the $T_a$ drop out (p1-29, p1-30).
\end{lemma}

\begin{proof}
Both sides are bilinear in $(X,Y)$ and invariant under $K$, because $\Ad(g)\Pi_P=\Pi_{\rho(g)P}$, $p_{RP}=Rp_PR^{\!\top}$ and $\rho$ is equivariant. So it suffices to take
$P=A$. Write $X=X_0+\sum x_{aw}\tfrac12aw$ with $X_0$ in the stabiliser; then $[X,\Pi_A]=\sum x_{aw}T_{aw}$, and the left side is $128\sum x_{aw}y_{aw}$. On the right,
$[\rho(\tfrac12aw),p_A]=-(E_{aw}+E_{wa})$, and $\rho_+X_0$ commutes with $p_A$, so the right side is $64\cdot2\sum x_{aw}y_{aw}$. Every tangent vector at $\Pi_P$ is $[Y,\Pi_P]$ for
some $Y\in\mathfrak k$, which gives the kinetic and current forms.
\end{proof}

The factor $64=512/8$ is the real trace on $\R^{512}$ against that on $\R^8$. At $n=2$ the same factor is $8=64/8$. With $\Tr(T_iT_j)=16\delta_{ij}$ there, this reproduces
the $8f^2$ and $16f^2$ of \cite[Thm.~26.7]{ES17}.

\begin{theorem}[Kinetic term and current at $\Pi_A$; fixed]
\label{p1:thm:kinetic}
With $\Pi=\Ad(e^X)\Pi_A$, $X=\sum x_{aw}\tfrac12aw$, and the conventions of \cite[Thm.~26.7]{ES17}, to quadratic order in $x$:
\[
\mathcal L_\Pi=64f^2\sum_{a,w}D_\mu x_{aw}D^\mu x_{aw},\qquad
\begin{aligned}
j^\mu_{\Pi,A}&=128f^2\sum_{w\in W}\Im(\bar\phi_wD^\mu\phi_w),\\
j^\mu_{\Pi,\rm em}&=128f^2\sum_{a\in A}\Im(\bar\phi_aD^\mu\phi_a),
\end{aligned}
\]
with $\phi_w=x_{2,w}+ix_{3,w}$ and $\phi_a=x_{a,13}+ix_{a,7}$. Since $[Q,\Pi_A]=0$ for both charges, no $A_\mu A^\mu$ term arises at $x=0$, and the criterion of
\cite[Cor.~26.2]{ES17} holds at $\Pi_A$ and $\Pi_W$ for $\uu1_A$ and for $\uu1_{\rm em}$. On the Clifford spheres of \cref{p1:prop:spheres}(iii) the formulas are exact,
with $\mathcal L_\Pi=64f^2\,Dn\cdot Dn$ (p1-31, p1-32).
\end{theorem}

\begin{proof}
$D\Pi=[DX,\Pi_A]+O(x^2)$, and $j_\Pi$ at second order is $-f^2\Tr([[Q,X],\Pi_A][DX,\Pi_A])$. p1-31 evaluates both bilinear forms with $\Tr(T_{aw}T_{a'w'})=128\delta$ and the
pairs of \cref{p1:thm:tangent}. The sign pattern is that of \cite[Thm.~26.7]{ES17}: $[Q,T_{\rm re}]=-T_{\rm im}$ and $[Q,T_{\rm im}]=T_{\rm re}$.
\end{proof}

\begin{theorem}[Curved extension; fixed]
\label{p1:thm:curved}
\begin{enumerate}[label=(\roman*),leftmargin=2.2em]
\item Let $\Omega=\Omega_{\rm Lorentz}+A^aT_a$ with $T_a$ spanning a subalgebra $\mathfrak s\subset\mathfrak k$. Then $D\Pi_P=0$ and $[\mathcal R,\Pi_P]=0$ hold for arbitrary spacetime and
gauge curvature iff $\mathfrak s\subset\so P\oplus\so{P^\perp\cap V_+}\oplus\so6_C$. At $\Pi_A$ and $\Pi_W$ this includes $\uu1_A$, $\uu1_{\rm em}$, $\mathfrak h_{\rm PS}$ and
$\mathfrak h_{\rm SM}$ (p1-35).
\item At every point of $F_3$, if $D_\mu\Pi=[Y_\mu,\Pi]$ with $Y_\mu\in\mathfrak k$, the adapted connection of \cite[Prop.~25.2]{ES17} is
$\mathcal C_\mu=[\Pi,D_\mu\Pi]=-\tfrac12(Y_\mu-\mathrm{vol}_PY_\mu\mathrm{vol}_P)$. This is a compact bivector of $\Lambda^2V_+$, so the adapted connection is spin-valued across
the whole Grassmannian family (p1-33).
\end{enumerate}
\end{theorem}

\begin{proof}
(i) The six interior bivectors commute with every $\Pi_P$ and with $\mathfrak k$. So $D\Pi_P=A^a[T_a,\Pi_P]$ and $\mathcal R=\mathcal R_{\rm Lorentz}+F^aT_a$, and both vanish
for arbitrary $A$ and $F$ iff each $T_a$ is in the stabiliser of \cref{p1:thm:classify}. (ii) With $v=\mathrm{vol}_P$ and $v^2=\id$,
$[\Pi,[Y,\Pi]]=\tfrac14[v,[Y,v]]=\tfrac12(vYv-Y)$. $\Ad(v)$ acts on vectors as $-1$ on $P$ and $+1$ on $P^\perp\cap V_+$ and $V_-$. So it maps $\Lambda^2V_+$ and
$\Lambda^2V_-$ to themselves, and the $\Lambda^2V_-$ part of $Y$ cancels.
\end{proof}

\begin{proposition}[Invariant couplings are constant; fixed]
\label{p1:prop:K}
Every $K$-invariant function on $F_3$ is constant. In particular the full alignment energy satisfies $\tfrac1{512}\sum\|[e_ie_j,\Pi]\|^2=16$ on $F_3$, with the sum over the $43$
compact blades. Every spectrum-only potential is constant on $F_3$. For a subalgebra $\mathfrak h\subsetneq\mathfrak k$ and positive couplings, the zero set of $V_{\mathfrak h}$ on
$F_3$ is the fixed set $\operatorname{Fix}(\mathfrak h)=\{\Pi\in F_3:[\mathfrak h,\Pi]=0\}$ (p1-34, p1-42).
\end{proposition}

\begin{proof}
$F_3$ is one $K$-orbit. The $43$ blades are orthogonal of equal norm in the trace form, which $\Ad(K)$ preserves, so $\sum_a\|[T_a,\Ad(g)\Pi]\|^2=\sum_a\|[\Ad(g^{-1})T_a,\Pi]\|^2$
is a sum over another orthonormal basis and has the same value. The value $16$ is sixteen movers, each with $\|[e_ie_j,\Pi_A]\|^2=512$. $V_{\mathfrak h}$ is a sum of
squares with positive weights.
\end{proof}

A coupling that selects a point of $F_3$ must therefore break $K$, and the natural candidates are the $V_{\mathfrak h}$ of proper subalgebras
\cite{Peskin80,Preskill81}. \Cref{p1:prop:duality} adds that every such $V_{\mathfrak h}$ is also invariant under $D$ and under $\Pi\mapsto\id-\Pi$.

\subsection{Fixed sets and the colour condition}
\label{p1:sec:fixed}

Write $U_A=\operatorname{span}(E_2,E_3)$, $U_{\rm em}=\operatorname{span}(w_1,w_2)$ and $W_0=\operatorname{span}(w_3,w_4)$.

\begin{theorem}[Fixed sets in $F_3$; fixed]
\label{p1:thm:fixed}
\begin{center}
\small
\begin{tabular}{@{}lllc@{}}
\toprule
algebra & fixed set in $F_3$ & dimensions & contains\\
\midrule
$\so6_C$, $\su3_C$ & all of $F_3$ & $16$ & all points\\
$\uu1_A$ & $\{P\supset U_A\}\sqcup\{P\perp U_A\}\cong\widetilde{\mathrm{Gr}}_2(\R^6)\sqcup\widetilde{\mathrm{Gr}}_4(\R^6)$ & $8,\ 8$ & $\Pi_{\pm A}$; $\Pi_{\pm W}$\\
$\uu1_{\rm em}$ & $\{P\supset U_{\rm em}\}\sqcup\{P\perp U_{\rm em}\}\cong\widetilde{\mathrm{Gr}}_2(\R^6)\sqcup\widetilde{\mathrm{Gr}}_4(\R^6)$ & $8,\ 8$ & $\Pi_{\pm W}$; $\Pi_{\pm A}$\\
$\R t_Y$ & $\{\Pi_{\pm A}\}\sqcup\{\Pi_{\pm W}\}\sqcup\{A'\oplus W'\}$ & $0,\ 0,\ 6$ & \\
$\su2_L$, $\mathfrak h_{\rm SM}$, $\mathfrak h_{\rm PS}$ & $\{\Pi_A,\Pi_{-A},\Pi_W,\Pi_{-W}\}$ & $0$ & four isolated points\\
$\mathfrak k$ & $\emptyset$ & & \\
\bottomrule
\end{tabular}
\end{center}
In the $t_Y$ row, $A'$ is a two-plane in $A$ and $W'$ a $J_Y$-complex line in $W$. That component is $\widetilde{\mathrm{Gr}}_2(A)\times\C P^1$, connected (p1-36 to p1-39).
\end{theorem}

\begin{proof}
For a connected group, $\Pi_P$ is fixed iff $P$ is invariant under $\rho_+(\mathfrak h)$ (\cref{p1:lem:transpose}). Let $j$ be a complex structure on a subspace $U$ that kills
$U^\perp$. If $P$ is $j$-invariant and $v=v_U+v_\perp\in P$, then $j^2v=-v_U\in P$, so $P=(P\cap U)\oplus(P\cap U^\perp)$ with $P\cap U$ itself $j$-invariant. For $U=U_A$ or
$U_{\rm em}$, $P\cap U\in\{0,U\}$, which gives the two Grassmannians of $U^\perp\cong\R^6$, each connected of dimension $8$. For $U=W$ and $j=J_Y$, $P\cap W$ is a complex
subspace of dimension $0$, $2$ or $4$. For $\su2_L$: $A$ is trivial and $W$ is irreducible without invariant vectors (p1-37). So $P\cap W\in\{0,W\}$, and if $P\cap W=0$
the $W$-component of every vector of $P$ is invariant, hence zero; thus $P\in\{A,W\}$. $\mathfrak h_{\rm SM}$ and $\mathfrak h_{\rm PS}$ contain $\su2_L$ and fix $A$ and $W$. $\so8$
is irreducible on $V_+$. The tangent kernels in p1-36, p1-38 and p1-39 agree with these dimensions, and p1-38 shows that the four points are isolated.
\end{proof}

\begin{proposition}[Unbroken algebras across $F_3$; fixed]
\label{p1:prop:strata}
The stabiliser of $\Pi_P$ in $\mathfrak h_{\rm PS}$ is $\so6_C\oplus\so{P\cap W}\oplus\so{P^\perp\cap W}$. In $\mathfrak h_{\rm SM}$ it is
$\su3_C\oplus\bigl(\mathfrak u_W\cap(\so{P\cap W}\oplus\so{P^\perp\cap W})\bigr)$, with $\mathfrak u_W=\rho(\su2_L\oplus\R t_Y)\cong\uu2$. At random points the dimensions in
$(\mathfrak h_{\rm SM},\mathfrak h_{\rm PS})$ are $(9,16)$ on both $\uu1_{\rm em}$-fixed components, namely $\su3_C\oplus\uu1_{\rm em}$ and $\so6_C\oplus\uu1_{\rm em}$. They are
$(8,16)$ on both $\uu1_A$-fixed components, $(8,15)$ at generic points of $F_3$ and $(12,21)$ at $\Pi_{\pm A},\Pi_{\pm W}$. On the six-dimensional
$t_Y$-component the stabiliser in $\mathfrak h_{\rm SM}$ has dimension $10$ (p1-39, p1-40).
\end{proposition}

\begin{proof}
If $X\in\so W$ preserves $P$ and kills $A$, then $X(P)\subset P\cap W$ and $X(P^\perp)\subset P^\perp\cap W$. So $X$ has image in $(P\cap W)\oplus(P^\perp\cap W)$, and by
antisymmetry it lies in $\so{P\cap W}\oplus\so{P^\perp\cap W}$; the converse is immediate. On $\{P\perp U_{\rm em}\}$ one generically has $P\cap W=0$ and
$P^\perp\cap W=U_{\rm em}$, and $\so{U_{\rm em}}=\R\rho(\tau_3)=\R\rho(B_{\rm em})\subset\mathfrak u_W$. p1-40 computes the dimensions of the intersections and of the stabilisers.
\end{proof}

The $\uu1_{\rm em}$-fixed component through $\Pi_A$ is $\widetilde{\mathrm{Gr}}_4(A\oplus W_0)$. Its tangent at $\Pi_A$ is $A\otimes W_0$, the neutral lower components of the four
doublets of \cref{p1:thm:tangent}. Along it $\su3\oplus\su2_L\oplus\uu1_Y$ breaks to $\su3\oplus\uu1_{\rm em}$, the pattern of a neutral doublet value (interpretive).

\begin{corollary}[Critical points and Hessians; fixed]
\label{p1:cor:critical}
Every smooth function on $F_3$ invariant under the connected group of $\mathfrak h\in\{\su2_L,\mathfrak h_{\rm SM},\mathfrak h_{\rm PS}\}$ has $\Pi_{\pm A}$ and $\Pi_{\pm W}$ as
critical points. Its Hessian there lies in the space of invariant symmetric forms on $A\otimes W$. For $\mathfrak h_{\rm PS}$ this space has dimension $10$, the forms
$S\otimes\delta_W$ with $S$ symmetric on $A$. For $\mathfrak h_{\rm SM}$ it has dimension $16$, the forms $S\otimes\delta_W+N\otimes\omega_Y$ with $N$ antisymmetric and
$\omega_Y(w,w')=\langle w,J_Yw'\rangle$ (p1-41).
\end{corollary}

\begin{proof}
At a fixed point the differential is an invariant covector, and by p1-38 there is none other than zero \cite{Michel71,Palais79}. The Hessian at a critical point is invariant. The
counts follow from Schur's lemma, with the commutant $\R$ of $\so4_W$ on $W$ and the commutant $\C=\operatorname{span}(\id,J_Y)$ of $\mathfrak u_W$; p1-41 computes them.
\end{proof}

\begin{remark}[The fourth rung, in dimensions; fixed]
\label{p1:rem:n4}
At $n=4$ the internal signature is $(11,9)$. The stabiliser of $E_2E_3E_8E_9$ in $\so{11,9}$ has dimension $126$ with compact part $63$, the $n=4$ row of
\cite[Prop.~26.9]{ES17}. The family $F_4=\widetilde{\mathrm{Gr}}_4(\R^{11})$ has dimension $91-63=28$ (p1-43). All $36$ bivectors of $V_-$ commute with every point, so any
colour algebra in $\so{V_-}$ passes the colour test everywhere. At $\Pi_A$ a $\uu1$ on a plane inside $A$ fixes a component of dimension $14=\dim\widetilde{\mathrm{Gr}}_2(\R^9)$,
and one on a plane orthogonal to $A$ fixes one of dimension $20=\dim\widetilde{\mathrm{Gr}}_4(\R^9)$. $\so4_W$ fixes a component of dimension $12=\dim\widetilde{\mathrm{Gr}}_4(\R^7)$
through $\Pi_A$, so at $n=4$ the Pati--Salam fixed points are no longer isolated (p1-44). In general $\dim F_n=4(3n-5)$, and a $\uu1$ on one positive plane fixes components of
dimensions $2(3n-5)$ and $4(3n-7)$: $(2)$, $(8,8)$, $(14,20)$ at $n=2,3,4$ (p1-45).
\end{remark}

\subsection{The answer}
\label{p1:sec:answer}

\begin{enumerate}[label=(\alph*),leftmargin=2.2em]
\item \emph{The family.} The $n\ge3$ analogue of the compact family is $F_n=\{\Pi_P:P\in\widetilde{\mathrm{Gr}}_4(V_+)\}$, the $K_n$-orbit of $\Pi_A$; at $n=2$ it is the paper's $S^4$
(fixed). At $n=3$ it is $\widetilde{\mathrm{Gr}}_4(\R^8)$, of dimension $16$, a compact symmetric space \cite{Helgason}. Its metric is $64$ times the trace metric of the
Grassmannian (\cref{p1:lem:vector}). The paper's $S^4$ is one of the inner Clifford spheres in it, and is not $K$-invariant (fixed).
\item \emph{Colour.} Every point of $F_3$ commutes with all of $\so6_C\supset\su3_C$, and so does every point of $F_n$ with $\so{V_-}$. The other involutive four-blade types never
commute with $\su3_C$ (fixed).
\item \emph{The vacuum.} $\Pi_A$ is the analogue of \cite[Def.~26.5]{ES17}. It is Lorentz invariant, neutral under $\uu1_A$ and $\uu1_{\rm em}$, fixed by $\mathfrak h_{\rm PS}$ and
$\mathfrak h_{\rm SM}$, moved by $16$ compact bivectors, and parallel under any connection valued in its stabiliser. It is the point $n=c$ of the embedded second-rung
family. The literal point $n=e$ is not neutral at $n=3$ (fixed).
\item \emph{Data for a selecting coupling.} A coupling that distinguishes the points of $F_3$ is not $K$-invariant (obstructed for $K$-invariant couplings). Its zero set, if it is
an alignment energy, is one of the fixed sets of \cref{p1:thm:fixed}. Its Hessian at the four Pati--Salam points lies in the $10$- or $16$-dimensional spaces of
\cref{p1:cor:critical}. Its fluctuations there are four doublets $(\mathbf1,\mathbf2)_{1/2}$ with a family index (fixed). No alignment energy of a subalgebra of $\mathfrak k$
separates $\Pi_A$, $\Pi_{-A}$, $\Pi_W$ and $\Pi_{-W}$; a coupling to the chiral block does, since on the block $D$ exchanges $P_L$ and $P_R$ (fixed). This is the
counterpart on $F_3$ of \cite[Prop.~21.6]{ES17} (interpretive).
\item \emph{The charge.} On $P_*$ the construction's charge $\tfrac12J_1$ is $\tfrac12J$, universal, and anomalous together with $\su2_L$ (fixed; restated). The electromagnetic
charge $B_{\rm em}$ also leaves $\Pi_A$ massless and has four charged complex fluctuations. Which charge, which coupling and which scale select the vacuum is not decided
here (open).
\end{enumerate}

\begin{table}[t]
\centering
\footnotesize
\begin{tabular}{@{}>{\raggedright\arraybackslash}p{8.4cm}>{\raggedright\arraybackslash}p{4.3cm}>{\raggedright\arraybackslash}p{2.2cm}@{}}
\toprule
statement & class & checks\\
\midrule
\cref{p1:prop:n2}: second-rung data, $K_2$-orbit $=S^4=\widetilde{\mathrm{Gr}}_4(\R^5)$, $Q$-fixed $S^2$, constant potentials & restated (ES17 Thms~26.6--26.8, Prop.~26.9); orbit and fixed set fixed & p1-1 to p1-7\\
\cref{p1:lem:census}: four-blade census, even-grade rule, $\mathrm{vol}_+=-Z$ & fixed & p1-8 to p1-10\\
\cref{p1:lem:transpose}: $\mathrm{vol}_P^{\!\top}=\mathrm{vol}_{\theta P}$, recovery of $P$ & fixed & p1-13\\
\cref{p1:lem:colour}: $\su3_C$ irreducible on $V_-$ & fixed & p1-15\\
\cref{p1:thm:classify}: three $K$-orbits, stabilisers, spectra, colour test & fixed; $n=3$ row of ES17 Prop.~26.9 restated & p1-11, p1-12, p1-14, p1-16, p1-42\\
\cref{p1:prop:spheres}: Clifford spheres, embedded second-rung family & fixed & p1-19, p1-32, p1-47\\
\cref{p1:prop:duality}: duality $P\mapsto P^\perp$, invariance of mass matrices & fixed & p1-20, p1-46\\
\cref{p1:thm:vacuum}: $\Pi_A$, $\Pi_W$ neutral, $\mathfrak h_{\rm PS}$- and $\mathfrak h_{\rm SM}$-fixed, not rigid; literal embedding fails & fixed & p1-21 to p1-23, p1-35\\
\cref{p1:thm:tangent}: fluctuations, four doublets $(\mathbf1,\mathbf2)_{1/2}$, charges & fixed & p1-24 to p1-26\\
\cref{p1:prop:block}: $\Pi_AP_*=P_L$, $q_A=J$ on the block, $P_*TP_*=0$, anomaly & fixed; (iv) restated (ES17 Prop.~8.5); anomaly condition inherited & p1-17, p1-18, p1-27, p1-28\\
\cref{p1:lem:vector}: vector model, kinetic term, current, $V_{\mathfrak h}$ & fixed & p1-29, p1-30\\
\cref{p1:thm:kinetic}: $\mathcal L_\Pi$ and $j_\Pi$ at $\Pi_A$ and on spheres & fixed & p1-31, p1-32\\
\cref{p1:thm:curved}: parallel vacua, compact adapted connection & fixed & p1-33, p1-35\\
\cref{p1:prop:K}: $K$-invariant couplings constant; zero sets & fixed; selection by $K$-invariant couplings obstructed & p1-34, p1-42\\
\cref{p1:thm:fixed}, \cref{p1:prop:strata}: fixed sets, unbroken algebras & fixed & p1-36 to p1-40\\
\cref{p1:cor:critical}: critical points and Hessians & fixed; criticality principle inherited & p1-38, p1-41\\
\cref{p1:rem:n4}: fourth rung & fixed; $n=4$ row of ES17 Prop.~26.9 restated & p1-43 to p1-45\\
family reading of $\so4_A$; electroweak pattern on the neutral component & interpretive & ---\\
which charge and coupling select the vacuum & open & ---\\
\bottomrule
\end{tabular}
\caption{Claims of \cref{p1:sec:main} and their classes. The checks are the \texttt{PASS} lines printed by \texttt{verify/p1\_families.py}.}
\label{p1:tab:claims}
\end{table}
